\par\noindent\textbf{(1)} 令ξ=x−st，行波方程−sf'+ff'=νf''。积分得νf'=[f²−2sf+C]/2，其中C待由两端f'→0确定。由于两远场必须同时是该一阶方程的根，此连接解要求代数多项式在uL和uR都为零。
\par\noindent\textbf{(2)} 把两端代入方程并相减，得uL²−2suL=uR²−2suR，故s=(uL+uR)/2，C=uLuR。由此νf'=(f−uL)(f−uR)/2，当uL>uR时两端之间f'<0，剖面随ξ单调下降。
\par\noindent\textbf{(3)} 分离变量得到f(ξ)=s−(Δu/2)tanh[Δu(ξ−ξ0)/(4ν)]，Δu=uL−uR>0。ξ→−∞回到uL，ξ→+∞回到uR，特征过渡厚度约4ν/Δu，与黏性扩散和平流非线性相平衡。
\par\noindent\textbf{(4)} ν→0+时剖面收缩成跳跃，速度满足通量F(u)=u²/2的Rankine–Hugoniot关系且uL>s>uR，符合Lax压缩熵条件。若uL<uR，特征线发散形成稀疏波，不应把上述tanh解机械反向当成熵允许激波。
